\chapter{مبرهنة ديريشليه للأوليات في المتتاليات الحسابية}

\section{نطاق الفصل وحالته}

\noindent\textbf{حالة الفصل:} \textenglish{REVIEWED}. يثبت هذا الفصل أن كل
فئة حسابية مختزلة تحتوي عددًا لا نهائيًا من الأعداد الأولية. والنتيجة
الوسيطة الأقوى هي تباعد مجموع مقلوبات الأوليات في تلك الفئة.

اجتازت النواة البرهانية فحوص الجودة والبناء، وتدقيقين منطقيين داخليين،
ومراجعة تحريرية ومرجعية، ثم مراجعة مستقلة عميقة وتحقيقًا مستقلاً لمواضع
المراجع. ولا تعني حالة \textenglish{REVIEWED} تحكيم مجلة أو بلوغ
\textenglish{RELEASE-READY}.

لا يثبت الفصل مبرهنة الأعداد الأولية في المتتاليات الحسابية، ولا يقدم
تقديرًا كميًا لخطأ التوزيع. المسار المستعمل نوعي، ويعتمد على تعامد
الشخصيات، والمنتج الأويلري، وعدم انعدام \(L(1,\chi)\) للشخصيات غير
الرئيسية، وقد ثبتت هذه الأدوات في الفصل السابع من دون اعتماد دائري.

\begin{remark}[مواضع المراجع]
قورنت بنية البرهان مع المسارات القياسية في
\cite[الفصل 7، الأقسام 7.3--7.9، الصفحات 148--155]{Apostol1976}،
و\cite[الفصلان 1 و4، الصفحات 1--11 و27--34]{Davenport2000}،
و\cite[الفصل 4، القسمان 4.2--4.3، الصفحات 115--132]{MontgomeryVaughan2007}.
وتستعمل لغة \cite[الفصول 3 و5 و17، الصفحات 43--64 و93--152 و419--426]{IwaniecKowalski2004}
لضبط الفصل بين النتيجة النوعية والنتائج الكمية اللاحقة.
\end{remark}

\section{مرشح الفئة الحسابية}

ليكن \(q\geq1\)، ولتكن \(a\) فئة مختزلة بترديد \(q\)، أي
\((a,q)=1\). نكتب
\[
P_{q,a}(\sigma)
=
\sum_{\substack{p\ \text{أولي}\\ p\equiv a\, (\mathrm{mod}\,q)}}
\frac1{p^\sigma},
\qquad \sigma>1.
\]

\begin{proposition}[مرشح الفئة بالشخصيات]
\label{prop:dirichlet-prime-class-filter}
\resultid{ANT-PROP-08-01}
\provedhere
لكل عدد صحيح \(n\):
\[
\mathbf 1_{n\equiv a\, (\mathrm{mod}\,q)}
=
\frac1{\varphi(q)}
\sum_{\chi\, (\mathrm{mod}\,q)}
\overline{\chi(a)}\chi(n).
\]
ومن ثم، لكل \(\sigma>1\):
\[
P_{q,a}(\sigma)
=
\frac1{\varphi(q)}
\sum_{\chi\, (\mathrm{mod}\,q)}
\overline{\chi(a)}
\sum_p\frac{\chi(p)}{p^\sigma}.
\]
\end{proposition}

\begin{proof}
إذا كان \((n,q)>1\)، فإن الطرف الأيمن يساوي صفرًا لأن كل شخصية ديريشليه
تنعدم عند \(n\)، والطرف الأيسر يساوي صفرًا أيضًا لأن \(a\) مختزلة.
أما إذا كان \((n,q)=1\)، فتعطي علاقة التعامد من الفصل السابع
\[
\sum_{\chi\, (\mathrm{mod}\,q)}
\chi(n)\overline{\chi(a)}
=
\begin{cases}
\varphi(q),&n\equiv a\pmod q,\\
0,&n\not\equiv a\pmod q.
\end{cases}
\]
وهذا يثبت الهوية الأولى. والهوية الثانية تنتج بضربها في \(p^{-\sigma}\)
والجمع على الأوليات. والتقارب المطلق يجيز تبديل الجمعين لأن عدد الشخصيات
منته، ولأن
\[
\sum_p\frac{|\chi(p)|}{p^\sigma}
\leq
\sum_{n=2}^{\infty}\frac1{n^\sigma}<\infty.
\]
\end{proof}

\section{لوغاريتم المنتج الأويلري}

عندما \(\Re(s)>1\)، نعرّف الفرع الأويلري للوغاريتم بالمتسلسلة
\[
\ell_\chi(s)
=
\sum_p\sum_{k=1}^{\infty}
\frac{\chi(p)^k}{k p^{ks}}.
\]
تتقارب هذه المتسلسلة مطلقًا ومحليًا بانتظام، وتحقق
\[
\exp\bigl(\ell_\chi(s)\bigr)=L(s,\chi).
\]
وبفصل الحد \(k=1\) نحصل على
\[
\ell_\chi(s)
=
\sum_p\frac{\chi(p)}{p^s}+H_\chi(s),
\]
حيث
\[
H_\chi(s)
=
\sum_p\sum_{k=2}^{\infty}
\frac{\chi(p)^k}{k p^{ks}}.
\]

\begin{lemma}[ضبط القوى الأولية العليا]
\label{lem:higher-prime-powers-remainder}
\resultid{ANT-LEM-08-01}
\provedhere
تتقارب سلسلة \(H_\chi(s)\) طبيعيًا على كل مجموعة مدمجة في نصف المستوى
\(\Re(s)>1/2\). وبخاصة، فهي هولومورفية ومحدودة في قرص مغلق مناسب مركزه
\(s=1\).
\end{lemma}

\begin{proof}
خذ \(0<r<1/2\)، وضع \(\sigma_0=1-r>1/2\). إذا كان
\(|s-1|\leq r\)، فإن \(\Re(s)\geq\sigma_0\)، ومن ثم
\begin{align*}
\sum_p\sum_{k=2}^{\infty}
\left|\frac{\chi(p)^k}{k p^{ks}}\right|
&\leq
\sum_p\sum_{k=2}^{\infty}\frac1{p^{k\sigma_0}}\\
&=
\sum_p\frac{p^{-2\sigma_0}}{1-p^{-\sigma_0}}\\
&\leq
\frac1{1-2^{-\sigma_0}}
\sum_{n=2}^{\infty}\frac1{n^{2\sigma_0}}<\infty.
\end{align*}
إذن التقارب طبيعي على القرص، وبمبرهنة فايرشتراس تكون الدالة هولومورفية
في داخله ومستمرة، ومن ثم محدودة، على القرص المغلق. والحجة نفسها تعمل على
كل مجموعة مدمجة في \(\Re(s)>1/2\).
\end{proof}

\section{فصل الشخصية الرئيسية}

\begin{proposition}[مساهمة الشخصية الرئيسية]
\label{prop:principal-character-log-growth}
\resultid{ANT-PROP-08-02}
\provedhere
إذا كانت \(\chi_0\) الشخصية الرئيسية بترديد \(q\)، فإن
\[
\ell_{\chi_0}(\sigma)
=
\log\frac1{\sigma-1}+O_q(1)
\qquad(\sigma\to1^+).
\]
\end{proposition}

\begin{proof}
للعدد الحقيقي \(\sigma>1\) تكون عوامل المنتج الأويلري موجبة، ولذلك
\(\ell_{\chi_0}(\sigma)=\log L(\sigma,\chi_0)\) باللوغاريتم الحقيقي.
ومن الفصل السابع:
\[
L(\sigma,\chi_0)
=
\zeta(\sigma)\prod_{p\mid q}(1-p^{-\sigma}).
\]
وبما أن
\[
(\sigma-1)\zeta(\sigma)\longrightarrow1,
\]
فإن
\[
\log\zeta(\sigma)
=
\log\frac1{\sigma-1}+o(1).
\]
أما المجموع المنتهي
\[
\sum_{p\mid q}\log(1-p^{-\sigma})
\]
فيتقارب إلى قيمة منتهية عندما \(\sigma\to1^+\). بضم الحدين نحصل على
الصيغة المطلوبة.
\end{proof}

\begin{proposition}[مساهمة الشخصيات غير الرئيسية]
\label{prop:nonprincipal-character-log-bounded}
\resultid{ANT-PROP-08-03}
\provedhere
إذا كانت \(\chi\) شخصية غير رئيسية بترديد \(q\)، فإن
\[
\ell_\chi(\sigma)=O_{q,\chi}(1)
\qquad(\sigma\to1^+).
\]
\end{proposition}

\begin{proof}
من المبرهنة \ref{thm:dirichlet-l-nonvanishing-at-one} لدينا
\[
L(1,\chi)\ne0.
\]
ولأن \(L(s,\chi)\) هولومورفية قرب \(s=1\)، يوجد قرص \(U\) مركزه \(1\)
لا تنعدم فيه الدالة. ولأن القرص بسيط الاتصال، يوجد دالة هولومورفية
\(g_\chi\) على \(U\) تحقق
\[
\exp(g_\chi(s))=L(s,\chi).
\]
على كل مركبة متصلة من
\[
U\cap\{s:\Re(s)>1\}
\]
لدينا
\[
\exp\bigl(\ell_\chi(s)-g_\chi(s)\bigr)=1.
\]
إذن الفرق يأخذ قيمًا في المجموعة المنفصلة \(2\pi i\mathbb Z\)، وهو
مستمر، فيكون ثابتًا. وبما أن \(g_\chi\) محدودة على قطعة مدمجة صغيرة تصل
إلى \(s=1\)، فإن \(\ell_\chi(\sigma)\) تبقى محدودة عندما
\(\sigma\to1^+\).
\end{proof}

\section{تباعد مجموع مقلوبات الأوليات}

\begin{theorem}[مبرهنة ديريشليه للأوليات في المتتاليات الحسابية]
\label{thm:dirichlet-primes-arithmetic-progressions}
\resultid{ANT-THM-08-01}
\provedhere
إذا كان \(q\geq1\) و\((a,q)=1\)، فإن
\[
P_{q,a}(\sigma)
=
\frac1{\varphi(q)}
\log\frac1{\sigma-1}+O_q(1)
\qquad(\sigma\to1^+).
\]
وبالتالي
\[
\sum_{\substack{p\ \text{أولي}\\ p\equiv a\, (\mathrm{mod}\,q)}}
\frac1p=\infty.
\]
وعلى وجه الخصوص، توجد أعداد أولية لا نهائية تحقق
\[
p\equiv a\pmod q.
\]
\end{theorem}

\begin{proof}
من القضية \ref{prop:dirichlet-prime-class-filter} ومن تعريف
\(H_\chi\):
\begin{align*}
P_{q,a}(\sigma)
&=
\frac1{\varphi(q)}
\sum_{\chi\, (\mathrm{mod}\,q)}
\overline{\chi(a)}
\bigl(\ell_\chi(\sigma)-H_\chi(\sigma)\bigr).
\end{align*}
يفصل حد الشخصية الرئيسية عن بقية الحدود. ولأن \((a,q)=1\)، فإن
\(\overline{\chi_0(a)}=1\). بحسب القضية
\ref{prop:principal-character-log-growth} يساوي حدها
\[
\frac1{\varphi(q)}
\log\frac1{\sigma-1}+O_q(1).
\]
وبحسب القضية \ref{prop:nonprincipal-character-log-bounded} تبقى حدود
الشخصيات غير الرئيسية محدودة. كما تبقى جميع حدود \(H_\chi(\sigma)\)
محدودة بحسب اللمّة \ref{lem:higher-prime-powers-remainder}. وعدد الشخصيات
بترديد \(q\) منته، فتنتج الصيغة التقاربية المعلنة مع حد إجمالي
\(O_q(1)\).

يبقى إثبات تباعد مجموع المقلوبات. افترض عكس المطلوب، أي
\[
\sum_{\substack{p\ \text{أولي}\\ p\equiv a\, (\mathrm{mod}\,q)}}
\frac1p<\infty.
\]
عندما \(\sigma\downarrow1\)، تزداد الحدود \(p^{-\sigma}\) رتيبًا إلى
\(p^{-1}\). لذلك تعطي مبرهنة التقارب الرتيب
\[
\lim_{\sigma\to1^+}P_{q,a}(\sigma)
=
\sum_{\substack{p\ \text{أولي}\\ p\equiv a\, (\mathrm{mod}\,q)}}
\frac1p<\infty.
\]
لكن الصيغة الأولى تقول إن
\[
P_{q,a}(\sigma)\longrightarrow\infty,
\]
وهذا تناقض. إذن مجموع المقلوبات متباعد، ولا يمكن أن تكون الأوليات في
الفئة محدودة العدد.
\end{proof}

\begin{remark}[حدود النتيجة]
المبرهنة السابقة تثبت اللانهاية وتباعد مجموع المقلوبات، لكنها لا تثبت
\[
\pi(x;q,a)\sim\frac{\operatorname{Li}(x)}{\varphi(q)}.
\]
فالنتيجة الأخيرة تحتاج أدوات كمية إضافية، أهمها ضبط الأصفار قرب الخط
\(\Re(s)=1\) وتحويل المعلومات التحليلية إلى تقديرات عدّية.
\end{remark}

